% =====================================================================
%  Harald Høffding, A SHORT SURVEY OF THE HISTORY OF MODERN PHILOSOPHY.
%  English translation of: Kort Oversigt over den nyere Filosofis Historie.  Fjerde, paany omarbejdede Udgave.
%  Kjøbenhavn og Kristiania: Gyldendalske Boghandel, Nordisk Forlag, 1910.
%
%  GENERATED; DO NOT EDIT BY HAND.  Built by ebook-method/kort-oversigt/tr_parts.py build --write from
%  tr_template.tex and the parts in ebook-method/kort-oversigt/tr/ (translated 2026-10-17 from
%  texts/hoeffding/kort-oversigt/transcription.tex, following TRANSLATOR-BRIEF.md and GLOSSARY.md there).
%  Page numbers in the margin [N] are those of the Danish print.  The one footnote is marked *) as printed.
%  Title: A Short Survey of the History of Modern Philosophy (Hans, 2026-10-17); not to be confused with the longer
%  "A Brief History of Modern Philosophy" (Macmillan 1912).  Terms: ebook-method/nyere-filosofis-historie/GLOSSARY.md.
%  ODD SPOTS: the six translators listed 54 places in the Danish that looked wrong or unclear (tr/odd-*.tsv); all were
%  read at the image: 33 faults of the transcription, mended in the Danish (patches.py) and here (the first line of
%  p. 2, a lost paragraph break on p. 32, specks, doubled dashes, italic titles); 9 misprints of the print, kept and
%  marked in the Danish, translated here as the sense requires.  The rest were grammar, fact or period spelling.
% =====================================================================
\documentclass[12pt,a4paper,openany]{book}
\usepackage[utf8]{inputenc}
\usepackage[T1]{fontenc}
\usepackage{libertinus}
\usepackage{textalpha}
\usepackage[english]{babel}
\usepackage{setspace}
\usepackage[margin=1.2in]{geometry}
\usepackage{fancyhdr}
\setlength{\headheight}{28pt}
\addtolength{\topmargin}{-13.5pt}
\usepackage{hyperref}
\pagestyle{fancy}
\fancyhf{}
\fancyhead[LE,RO]{\thepage}
\fancyhead[RE]{\textit{Harald Høffding: A Short Survey}}
\fancyhead[LO]{\textit{\leftmark}}
\fancypagestyle{plain}{\fancyhf{}\fancyhead[LE,RO]{\thepage}}
\setstretch{1.3}
\usepackage{marginnote}
\usepackage{graphicx}
\DeclareUnicodeCharacter{0254}{\reflectbox{c}}  % ɔ: the old Danish id-est mark (a reversed c)
\newcommand{\starmark}[1]{\ifcase#1\or*\or**\or***\fi)}   % the print marks its one note *)
\renewcommand{\thefootnote}{\starmark{\value{footnote}}}
\newcommand{\opage}[1]{\marginnote{\footnotesize\textit{[#1]}}}
\newcommand{\apage}[1]{\marginnote{\footnotesize\textit{[#1]}}}
\newcommand{\whiteline}{\par\vspace{\baselineskip}\par}
% the printed Indhold: an essay or chapter line (title, page), a section line
\newcommand{\ientry}[2]{\par\smallskip\noindent #1\dotfill\,#2\par}
\newcommand{\isub}[2]{\par\noindent\hspace*{2em}#1\dotfill\,#2\par}
% the headings (no \addcontentsline: the print has its own Indhold, and a footnote in a heading breaks a toc entry)
\pdfstringdefDisableCommands{\renewcommand{\footnote}[2][]{}}
\newcommand{\prefacehead}[1]{\chapter*{#1}\markboth{Preface}{Preface}}
\newcommand{\sectionhead}[1]{\par\medskip\begin{center}\textbf{#1}\end{center}\smallskip}   % the print: centred, bold
\newcommand{\introhead}[1]{\begin{center}{\large #1}\\[4pt]\rule{2em}{0.4pt}\end{center}\markboth{Introduction}{Introduction}\medskip}
\newcommand{\closerule}{\par\bigskip\begin{center}\rule{4em}{0.4pt}\end{center}\par}   % the short rule that closes each part of the print
\newcommand{\chapterheadone}[2]{\bigskip\begin{center}{\large #1}\\[4pt]\rule{2em}{0.4pt}\end{center}\markboth{#2}{#2}\medskip}   % VIII: numeral and title on one line
\newcommand{\chapterhead}[2]{\bigskip\begin{center}#1\\[3pt]{\large #2}\\[4pt]\rule{2em}{0.4pt}\end{center}\markboth{#2}{#2}\medskip}   % numeral, title in capitals, short rule

\begin{document}

\frontmatter
\thispagestyle{empty}
\vspace*{5cm}
\begin{center}
{\large A SHORT SURVEY OF THE\\[4pt] HISTORY OF MODERN PHILOSOPHY}
\end{center}
\clearpage
\begin{titlepage}
\centering
\vspace*{2cm}
{HARALD HØFFDING}\par
\vspace{6pt}\rule{2em}{0.4pt}\par
\vspace{1.2cm}
{\Huge A SHORT SURVEY OF\\[6pt] THE HISTORY OF\\[6pt] MODERN PHILOSOPHY}\par
\vspace{1cm}
{\small FOURTH, NEWLY REVISED EDITION}\par
\vfill
{\small [the publisher's device]}\par
\vfill
{COPENHAGEN AND KRISTIANIA}\par
{GYLDENDALSKE BOGHANDEL}\par
{NORDISK FORLAG}\par
\vspace{4pt}\rule{2em}{0.4pt}\par\vspace{4pt}
{\footnotesize FORLAGETS BOGTRYKKERI}\par
{\footnotesize 1910}\par
\vspace{1cm}
{\small\textit{Translated from the Danish}}\par
\end{titlepage}
\clearpage

\chapter*{CONTENTS}
\markboth{Contents}{Contents}
\noindent\hfill{\small Page}\par
\ientry{Introduction. The Philosophical Problems}{1}
\ientry{I. The Origin of Modern Philosophy}{5}
\isub{1. The Discovery of Natural Man}{5}
\isub{2. The New Picture of the World}{11}
\isub{3. The New Science}{14}
\ientry{II. The Formation of Systems}{17}
\isub{a. Réné Descartes}{18}
\isub{b. Thomas Hobbes}{21}
\isub{c. Baruch Spinoza}{26}
\isub{d. Gottfried Wilhelm Leibniz}{32}
\ientry{III. The Philosophy of Experience}{35}
\isub{a. John Locke}{36}
\isub{b. George Berkeley}{38}
\isub{c. David Hume}{39}
\ientry{IV. The Philosophy of the Enlightenment}{43}
\isub{1. Voltaire, Holbach and Rousseau}{44}
\isub{2. Wolff and Lessing}{48}
\ientry{V. Immanuel Kant and the Critical Philosophy}{50}
\isub{Biography and Characterization}{50}
\isub{1. Epistemology}{53}
\isub{2. Ethics}{59}
\isub{3. Theory of Unity on an Aesthetic and Biological Basis}{61}
\ientry{VI. The Philosophy of Romanticism}{63}
\isub{a. Fichte, Schelling and Hegel}{63}
\isub{b. Schopenhauer, Kierkegaard and Feuerbach}{67}
\ientry{VII. Positivism}{74}
\isub{1. Saint Simon and Comte}{75}
\isub{2. Stuart Mill}{77}
\isub{3. Darvin and Spencer}{78}
\ientry{VIII. Recent Treatments of Problems}{81}
\isub{a. The Psychological Problem}{82}
\isub{b. The Problem of Existence}{85}
\isub{c. The Problem of Knowledge}{88}
\isub{d. The Problem of Valuation}{92}
\clearpage

\chapter*{PREFACE TO THE FOURTH EDITION}
\markboth{Preface}{Preface}

In this new edition I have entirely reworked the little book in the light of the
experience I have gained by using it for a number of years in the propaedeutic
course. I have left out much that was seldom needed in that course, but have in
return included biographical information, and I have also tried to make the
presentation more rounded and intelligible. This applies not least to the last
section, which was to give an idea of the present work of philosophy. --- I offer
my best thanks to Dr.\ Anton Thomsen for a number of valuable remarks, which have
benefited this new edition.

\textit{The 28th of March, 1910.}

\hfill\textbf{Harald Høffding.}
\clearpage

\mainmatter
%%PARTS%%

\end{document}
